
\documentclass[11pt,a4paper]{article}

\usepackage[utf8]{inputenc}
\usepackage[T1]{fontenc}
\usepackage[french]{babel}
\usepackage{amsmath,amssymb,amsfonts,amsthm,mathtools}
\usepackage{booktabs}
\usepackage{geometry}
\usepackage{graphicx}
\usepackage{enumitem}
\usepackage{hyperref}
\usepackage{algorithm}
\usepackage{algorithmic}

\geometry{margin=2.5cm}

% -------------------------------------------------
% Environnements
% -------------------------------------------------
\newtheorem{definition}{Définition}[section]
\newtheorem{assumption}{Hypothèse}[section]
\newtheorem{proposition}{Proposition}[section]
\newtheorem{theorem}{Théorème}[section]
\newtheorem{corollary}{Corollaire}[section]
\newtheorem{remark}{Remarque}[section]
\newtheorem{lemma}{Lemme}[section]

% -------------------------------------------------
% Opérateurs
% -------------------------------------------------
\DeclareMathOperator{\prox}{prox}
\DeclareMathOperator{\argmin}{argmin}
\DeclareMathOperator{\argmax}{argmax}
\DeclareMathOperator{\dist}{dist}
\DeclareMathOperator{\conv}{conv}
\DeclareMathOperator{\softmax}{softmax}

% -------------------------------------------------
% Raccourcis
% -------------------------------------------------
\newcommand{\R}{\mathbb R}
\newcommand{\I}{\mathrm{Id}}
\newcommand{\Mcal}{\mathcal M}
\newcommand{\Ccal}{\mathcal C}
\newcommand{\Pcal}{\mathcal P}
\newcommand{\Tcal}{\mathcal T}
\newcommand{\Vcal}{\mathcal V}
\newcommand{\Rcal}{\mathcal R}

\title{Routage variationnel pour PCH\\
\large Choix du contexte comme problème d'optimisation}
\author{}
\date{}

\begin{document}

\maketitle

\begin{abstract}
Ce document formalise le routage dans le cadre PCH. Le noyau PCH est supposé déjà défini et inchangé. L'objet de cette extension est de poser le choix du contexte comme un problème d'optimisation rigoureux sur l'espace produit \(\R^d\times\Ccal\), puis comme une relaxation convexe sur le simplexe des contextes. La quantité fondamentale de sélection est la valeur réduite normalisée associée au problème de routage. Le document introduit un routage dur exact, un routage souple entropique, un point de vue géométrique sur les cellules de contexte, et un certificat de pré-roulage court comme approximation calculable.
\end{abstract}

\tableofcontents

% =================================================
\section{Problème visé}
% =================================================

On considère un noyau PCH déjà défini et validé. Pour chaque entrée latente \(w\in\R^d\), il existe plusieurs contextes admissibles \(c\in\Ccal\), chacun muni d'un ensemble convexe admissible, fermé et non vide \(\Mcal_c\subset\R^d\) et d'une fonctionnelle locale propre.

Le problème n'est pas de redéfinir le noyau PCH. Le problème est de choisir \emph{le contexte retenu par le routage} avant la purification complète, de manière rigoureuse, stable et compatible avec la théorie.

Dans cette version, le noyau PCH reste défini par les énergies \(J_c\). Le routage constitue un problème auxiliaire de décision, défini séparément par une énergie normalisée. Cette distinction est volontaire : le noyau PCH et le routage ne résolvent pas le même problème d'optimisation.

Le noyau PCH est conçu pour produire une purification adaptée au contexte retenu. La difficulté principale porte donc sur la sélection de ce contexte.

% =================================================
\section{Cadre de base du noyau PCH}
% =================================================

\begin{assumption}[Espace latent]
On fixe un espace euclidien réel \(\R^d\). Les états latents purifiés vivent dans \(\R^d\).
\end{assumption}

\begin{assumption}[Famille de contextes]
On considère un ensemble fini de contextes admissibles
\[
\Ccal = \{1,\dots,K\}.
\]
À chaque contexte \(c\in\Ccal\) est associé un ensemble convexe non vide, fermé, \(\Mcal_c\subset\R^d\).
\end{assumption}

\begin{assumption}[Paramètres PCH]
Pour chaque contexte \(c\), on fixe des paramètres \(\lambda_{\mathrm{ctx},c}>0\), \(\lambda_{\mathrm{ridge},c}\ge 0\), et un paramètre proximal \(\eta_c>0\).
\end{assumption}

\begin{definition}[Énergie contextuelle]
Pour un contexte \(c\in\Ccal\), on définit l'énergie
\[
J_c(u;w)
=
\frac12\|u-w\|^2
+
\lambda_{\mathrm{ctx},c}\,\dist(u,\Mcal_c)^2
+
\lambda_{\mathrm{ridge},c}\,\|u\|^2.
\]
\end{definition}

\begin{definition}[Opérateur PCH associé]
Pour un contexte \(c\in\Ccal\), on définit
\[
\Pcal_c(w)
:=
\argmin_{u\in\R^d} J_c(u;w).
\]
\end{definition}

\begin{proposition}[Convexité forte et unicité]
Pour tout contexte \(c\) et toute entrée \(w\), la fonctionnelle \(u\mapsto J_c(u;w)\) est \(\mu_c\)-fortement convexe, où
\[
\mu_c:=1+2\lambda_{\mathrm{ridge},c}.
\]
En particulier, \(\Pcal_c(w)\) est un singleton.
\end{proposition}

\begin{proof}
Le terme \(\frac12\|u-w\|^2+\lambda_{\mathrm{ridge},c}\|u\|^2\) est \((1+2\lambda_{\mathrm{ridge},c})\)-fortement convexe en \(u\). La fonction \(u\mapsto \dist(u,\Mcal_c)^2\) est convexe puisque \(\Mcal_c\) est non vide, fermé et convexe. La somme d'une fonction fortement convexe et de fonctions convexes est fortement convexe. L'existence et l'unicité du minimiseur en découlent.
\end{proof}

\begin{remark}
Le noyau PCH n'est pas modifié par ce document. On ne change ni la forme de \(J_c\), ni la structure des ensembles \(\Mcal_c\), ni la purification finale. On ne traite ici que le routage.
\end{remark}

% =================================================
\section{Normalisation canonique, Moreau et valeur réduite}
% =================================================

La normalisation est choisie pour rendre explicite l'échelle naturelle de chaque contexte. Le facteur \(\mu_c=1+2\lambda_{\mathrm{ridge},c}\) est la constante de convexité forte de la partie quadratique de \(J_c\). Après division par \(\mu_c\), la fonctionnelle normalisée se lit comme une enveloppe de Moreau de paramètre \(\mu_c\), ce qui rend comparables des contextes ayant des niveaux de ridge différents et fournit une formule de gradient standard. Ce point est classique dans l'analyse convexe ; nous le rappelons pour rendre le manuscrit autonome \cite{Moreau1965,Rockafellar1970,BauschkeCombettes2017}.

On introduit, pour chaque contexte, les paramètres normalisés
\[
\mu_c = 1 + 2\lambda_{\mathrm{ridge},c},
\qquad
\gamma_c = \frac{\lambda_{\mathrm{ctx},c}}{\mu_c},
\qquad
\tau_c = \eta_c \mu_c.
\]

\begin{definition}[Fonction de base normalisée]
Pour tout contexte \(c\), on pose
\[
f_c(u):=\lambda_{\mathrm{ctx},c}\,\dist(u,\Mcal_c)^2+\lambda_{\mathrm{ridge},c}\,\|u\|^2,
\qquad
g_c(u):=\frac{1}{\mu_c}f_c(u).
\]
Alors
\[
\widetilde J_c(u;w)=\frac{1}{\mu_c}J_c(u;w)
=
g_c(u)+\frac{1}{2\mu_c}\|u-w\|^2.
\]
\end{definition}

\begin{proposition}[Justification de la normalisation]
Pour chaque contexte \(c\), la normalisation par \(\mu_c\) a les propriétés suivantes :
\begin{enumerate}[leftmargin=2em]
\item elle ne change pas le minimiseur de \(J_c(\cdot;w)\) ;
\item elle transforme le problème contextuel en un problème proximal de paramètre \(\mu_c\) pour la fonction convexe \(g_c\) ;
\item elle met la comparaison des valeurs réduites sur une échelle canonique liée à la convexité forte du terme quadratique ;
\item elle évite qu'un contexte à ridge élevé domine artificiellement la comparaison des valeurs optimales.
\end{enumerate}
\end{proposition}

\begin{proof}
Le point (1) est immédiat car la multiplication par un scalaire strictement positif ne change pas l'ensemble des minimisateurs. Le point (2) suit de l'identité
\[
\widetilde J_c(u;w)=g_c(u)+\frac{1}{2\mu_c}\|u-w\|^2,
\]
qui est précisément la forme standard d'une enveloppe de Moreau de paramètre \(\mu_c\). Le point (3) est une conséquence de la définition même de \(\mu_c\), qui est la constante de forte convexité de la partie quadratique en \(u\). Le point (4) formalise l'intérêt comparatif de la normalisation : deux contextes ayant des pénalités de ridge différentes sont alors évalués sur une échelle homogène.
\end{proof}

\begin{definition}[Valeur réduite du contexte]
Pour tout contexte \(c\) et toute entrée \(w\), on définit la valeur réduite
\[
\widetilde V_c(w)
:=
\inf_{u\in\R^d}\widetilde J_c(u;w).
\]
Par unicité du minimiseur, on note ce point
\[
u_c^\star(w):=\argmin_{u\in\R^d}\widetilde J_c(u;w).
\]
\end{definition}

\begin{definition}[Enveloppe de Moreau associée au contexte]
On note l'enveloppe de Moreau de \(g_c\) au paramètre \(\mu_c\) par
\[
e_{\mu_c g_c}(w)
:=
\inf_{u\in\R^d}
\left\{
g_c(u)+\frac{1}{2\mu_c}\|u-w\|^2
\right\}.
\]
\end{definition}

\begin{proposition}[Enveloppe de Moreau et gradient]
Pour tout contexte \(c\), la fonction \(g_c\) est propre, convexe et semi-continue inférieurement. On a
\[
u_c^\star(w)=\prox_{\mu_c g_c}(w),
\qquad
\widetilde V_c(w)=e_{\mu_c g_c}(w).
\]
De plus, \(e_{\mu_c g_c}\) est de classe \(C^1\) sur \(\R^d\) et
\[
\nabla e_{\mu_c g_c}(w)=\frac{1}{\mu_c}\bigl(w-\prox_{\mu_c g_c}(w)\bigr)
=\frac{1}{\mu_c}\bigl(w-u_c^\star(w)\bigr).
\]
Par conséquent,
\[
\nabla \widetilde V_c(w)=\frac{1}{\mu_c}\bigl(w-u_c^\star(w)\bigr).
\]
\end{proposition}

\begin{proof}
La convexité et la semi-continuité inférieure de \(g_c\) découlent de celles de \(\dist(\cdot,\Mcal_c)^2\) et de \(\|\cdot\|^2\). Le résultat classique sur les enveloppes de Moreau d'une fonction propre, convexe et s.c.i. implique que \(e_{\mu_c g_c}\) est de classe \(C^1\) et que
\[
\nabla e_{\mu_c g_c}(w)=\frac{1}{\mu_c}\bigl(w-\prox_{\mu_c g_c}(w)\bigr).
\]
Comme
\[
\widetilde J_c(u;w)=g_c(u)+\frac{1}{2\mu_c}\|u-w\|^2,
\]
les minimisateurs coïncident et la formule annoncée suit.
\end{proof}

\begin{remark}
Le gain normalisé
\[
\widetilde B_c(w):=\widetilde J_c(w;w),
\qquad
\widetilde\Gamma_c(w):=\frac{\widetilde B_c(w)-\widetilde V_c(w)}{\widetilde B_c(w)+\varepsilon},
\qquad \varepsilon>0,
\]
est un indicateur secondaire de diagnostic. La quantité fondamentale du routage est \(\widetilde V_c(w)\), non le gain.
\end{remark}

% =================================================
\section{Routage exact sur \(\R^d\times\Ccal\)}
% =================================================

\begin{definition}[Fonctionnelle de routage]
Pour une entrée \(w\in\R^d\), on introduit la fonctionnelle de routage
\[
\mathcal E_{\mathrm{route}}(w;u,c):=\widetilde J_c(u;w),
\qquad
(u,c)\in \R^d\times\Ccal.
\]
\end{definition}

\begin{definition}[Problème global de routage exact]
Le routage exact est le problème
\[
(\hat u(w),\hat c(w))
\in
\argmin_{(u,c)\in \R^d\times\Ccal}\mathcal E_{\mathrm{route}}(w;u,c).
\]
\end{definition}

\begin{proposition}[Équivalence avec la minimisation imbriquée]
Pour tout \(w\in\R^d\),
\[
\inf_{(u,c)\in \R^d\times\Ccal}\mathcal E_{\mathrm{route}}(w;u,c)
=
\min_{c\in\Ccal}\inf_{u\in\R^d}\widetilde J_c(u;w).
\]
De plus, si pour chaque \(c\) le minimiseur \(u_c^\star(w)\) existe, alors tout minimiseur global \((\hat u(w),\hat c(w))\) vérifie
\[
\hat u(w)=u_{\hat c(w)}^\star(w)
\quad\text{et}\quad
\hat c(w)\in\argmin_{c\in\Ccal}\,\widetilde V_c(w).
\]
\end{proposition}

\begin{proof}
Comme \(\Ccal\) est fini,
\[
\inf_{(u,c)\in \R^d\times\Ccal}\mathcal E_{\mathrm{route}}(w;u,c)
=
\min_{c\in\Ccal}\inf_{u\in\R^d}\widetilde J_c(u;w).
\]
Pour un contexte fixé \(c\), l'infimum en \(u\) est atteint au minimiseur \(u_c^\star(w)\). Le couple optimal s'obtient donc en choisissant un contexte minimisant la valeur réduite et en prenant le minimiseur correspondant.
\end{proof}

\begin{proposition}[Existence du minimiseur global]
Pour tout \(w\in\R^d\), l'ensemble
\[
\argmin_{(u,c)\in \R^d\times\Ccal}\mathcal E_{\mathrm{route}}(w;u,c)
\]
est non vide.
\end{proposition}

\begin{proof}
Pour chaque \(c\in\Ccal\), la valeur \(\widetilde V_c(w)\) est atteinte en un unique point \(u_c^\star(w)\). Comme \(\Ccal\) est fini, la fonction \(c\mapsto \widetilde V_c(w)\) atteint son minimum sur \(\Ccal\). Si \(c^\star\) est un minimiseur, alors \((u_{c^\star}^\star(w),c^\star)\) est un minimiseur global.
\end{proof}

% =================================================
\section{Théorème central de synthèse}
% =================================================

\begin{theorem}[Structure variationnelle globale du routage PCH]
Sous les hypothèses du document, le routage PCH admet la structure suivante :
\begin{enumerate}[leftmargin=2em]
\item Pour chaque contexte \(c\), le problème local admet un unique minimiseur
\[
u_c^\star(w)=\prox_{\mu_c g_c}(w),
\qquad
\widetilde V_c(w)=e_{\mu_c g_c}(w).
\]

\item La valeur globale du routage
\[
F(w)=\min_{c\in\Ccal}\widetilde V_c(w)
\]
est continue, localement lipschitzienne, et de classe \(C^1\) sur toute cellule de dominance stricte.

\item L'espace latent se décompose en cellules de dominance stricte \(\Omega_c^{\mathrm{str}}\), en cellules faibles \(\Omega_c\), et en une strate de bord \(\mathcal B\). Sur chaque cellule stricte, le contexte routé est localement constant.

\item Le routage dur exact est la résolution du problème mixte
\[
\min_{(u,c)\in\R^d\times\Ccal}\mathcal E_{\mathrm{route}}(w;u,c),
\]
et son écriture disjonctive équivalente est une programmation mixte à \(K\) branches.

\item Le routage souple est la régularisation entropique du critère des valeurs réduites. Son minimiseur unique est la loi de Gibbs
\[
p_c^\star(w)
=
\frac{\exp(-\beta\,\widetilde V_c(w))}
{\sum_{j=1}^K \exp(-\beta\,\widetilde V_j(w))},
\]
et
\[
p^\star_\beta(w)\xrightarrow[\beta\to+\infty]{}e_{c^\star(w)}
\]
lorsque le minimiseur dur \(c^\star(w)\) est unique.

\item Le changement de contexte est localement stable hors de la strate de bord, et les frontières de décision se déplacent continûment sous perturbation des paramètres lorsque la non-dégénérescence est satisfaite.

\item La fonction globale \(F\) est l'enveloppe inférieure de la famille finie d'enveloppes de Moreau \(\{e_{\mu_c g_c}\}_{c\in\Ccal}\).
\end{enumerate}
\end{theorem}

\begin{proof}
Les points \(1\), \(2\), \(3\), \(4\), \(5\), \(6\) et \(7\) résultent respectivement des sections précédentes sur la normalisation, la valeur globale, la partition de l'espace latent, la formulation mixte, le routage souple, la sensibilité des frontières et l'interprétation variationnelle complète.
\end{proof}

% =================================================
\section{Stabilité globale, cellules et frontières}
% =================================================

\begin{proposition}[Continuité et régularité locale de \(F\)]
La fonction \(F:\R^d\to\R\) est continue sur \(\R^d\) et localement lipschitzienne. Plus précisément, pour tout compact \(\mathcal K\subset\R^d\), \(F\) est lipschitzienne sur \(\mathcal K\).
\end{proposition}

\begin{proof}
Chaque fonction \(\widetilde V_c\) est de classe \(C^1\), donc continue. Comme \(\Ccal\) est fini, \(F\) est le minimum fini d'un nombre fini de fonctions continues, donc \(F\) est continue.

Soit \(\mathcal K\subset\R^d\) compact. Pour chaque \(c\in\Ccal\), la continuité de \(\nabla \widetilde V_c\) implique
\[
L_c:=\sup_{w\in\mathcal K}\|\nabla \widetilde V_c(w)\|<+\infty.
\]
Par le théorème des accroissements finis, \(\widetilde V_c\) est lipschitzienne sur \(\mathcal K\) avec constante \(L_c\). En posant \(L_{\mathcal K}:=\max_{c\in\Ccal}L_c\), on obtient
\[
|F(x)-F(y)|
\le
L_{\mathcal K}\,\|x-y\|
\qquad (x,y\in\mathcal K).
\]
Donc \(F\) est lipschitzienne sur \(\mathcal K\), d'où la localement lipschitzianité.
\end{proof}

\begin{proposition}[Correspondance active et stabilité locale]
La correspondance \(\mathcal A:\R^d\rightrightarrows \Ccal\) a des valeurs non vides et finies. De plus :
\begin{enumerate}[leftmargin=2em]
\item son graphe est fermé ;
\item elle est localement constante sur toute cellule de dominance stricte ;
\item elle est semi-continue supérieurement sur \(\R^d\).
\end{enumerate}
\end{proposition}

\begin{proof}
Pour chaque \(c\in\Ccal\),
\[
\{w\in\R^d:\ c\in\mathcal A(w)\}
=
\bigcap_{j\in\Ccal}\{w\in\R^d:\widetilde V_c(w)\le \widetilde V_j(w)\},
\]
intersection finie d'ensembles fermés, donc ensemble fermé. Cela prouve que le graphe de \(\mathcal A\) est fermé et donc que \(\mathcal A\) est semi-continue supérieurement, puisque \(\Ccal\) est fini et muni de la topologie discrète. Ce résultat peut aussi être vu comme une conséquence d'un théorème du maximum de Berge dans le cas discret \cite{Berge1963}.

Si \(w\in\Omega_c^{\mathrm{str}}\), les inégalités sont strictes et stables par continuité, donc \(c\) reste l'unique minimiseur dans un voisinage de \(w\). Cela donne la localité de la constance.
\end{proof}

\begin{definition}[Cellules de contexte]
On définit la cellule stricte associée à un contexte \(c\in\Ccal\) par
\[
\Omega_c^{\mathrm{str}}
:=
\left\{
w\in\R^d:\ \widetilde V_c(w)<\widetilde V_j(w)\ \text{pour tout }j\in\Ccal,\ j\neq c
\right\}.
\]
On définit également la cellule de dominance faible
\[
\Omega_c
=
\left\{
w\in\R^d:\ c\in\argmin_{j\in\Ccal}\widetilde V_j(w)
\right\}
=
\bigcap_{j\neq c}\{w\in\R^d:\widetilde V_c(w)\le \widetilde V_j(w)\}.
\]
Enfin, on note
\[
\mathcal B
:=
\left\{
w\in\R^d:\ \#\mathcal A(w)\ge 2
\right\}
\]
l'ensemble des points de bord, où au moins deux contextes sont ex æquo.
\end{definition}

\begin{proposition}[Partition de l'espace latent]
Pour chaque \(c\in\Ccal\), l'ensemble \(\Omega_c^{\mathrm{str}}\) est ouvert et les ensembles \(\Omega_c\) sont fermés. De plus, les cellules strictes sont deux à deux disjointes, et
\[
\R^d
=
\left(\bigsqcup_{c\in\Ccal}\Omega_c^{\mathrm{str}}\right)\sqcup \mathcal B.
\]
\end{proposition}

\begin{proof}
Chaque fonction \(\widetilde V_c-\widetilde V_j\) est continue. Donc \(\Omega_c^{\mathrm{str}}\), défini par des inégalités strictes de fonctions continues, est ouvert, et \(\Omega_c\), défini par des inégalités larges, est fermé.

Si \(w\in\Omega_c^{\mathrm{str}}\cap\Omega_j^{\mathrm{str}}\) avec \(c\neq j\), alors on aurait simultanément
\[
\widetilde V_c(w)<\widetilde V_j(w)
\quad\text{et}\quad
\widetilde V_j(w)<\widetilde V_c(w),
\]
absurde. Les cellules strictes sont donc disjointes.

Tout point \(w\in\R^d\) appartient soit à une cellule stricte, si le minimiseur est unique, soit à \(\mathcal B\), si au moins deux contextes sont ex æquo. Cela donne la décomposition annoncée.
\end{proof}

\begin{corollary}[Stabilité géométrique du contexte]
Si \(w\in\Omega_c^{\mathrm{str}}\), alors il existe un voisinage ouvert \(U\) de \(w\) tel que
\[
U\subset \Omega_c^{\mathrm{str}}.
\]
En particulier, le contexte routé est localement constant sur chaque cellule stricte.
\end{corollary}

\begin{proof}
C'est une conséquence immédiate de l'ouverture de \(\Omega_c^{\mathrm{str}}\).
\end{proof}

\begin{proposition}[Sensibilité locale des frontières]
Pour chaque paire \(c\neq j\), on pose
\[
\Gamma_{c,j}(w):=\widetilde V_c(w)-\widetilde V_j(w),
\qquad
\Sigma_{c,j}:=\{w\in\R^d:\Gamma_{c,j}(w)=0\}.
\]
Alors \(\Sigma_{c,j}\) est fermé. Si \(w_0\in\Sigma_{c,j}\) et si
\[
\nabla \Gamma_{c,j}(w_0)\neq 0,
\]
alors il existe un voisinage \(U\) de \(w_0\) dans lequel \(\Sigma_{c,j}\cap U\) est une hypersurface \(C^1\) de \(\R^d\).
\end{proposition}

\begin{proof}
La fermeture de \(\Sigma_{c,j}\) découle de la continuité de \(\Gamma_{c,j}\). Si \(\nabla \Gamma_{c,j}(w_0)\neq 0\), le théorème des fonctions implicites s'applique à l'équation \(\Gamma_{c,j}(w)=0\) au voisinage de \(w_0\), ce qui fournit une hypersurface \(C^1\) locale.
\end{proof}

\begin{proposition}[Stabilité sous marge stricte]
Supposons qu'en un point \(w\in\R^d\), le minimiseur \(c^\star\) de \(c\mapsto \widetilde V_c(w)\) soit unique et qu'il existe une marge
\[
\delta(w):=\min_{c\neq c^\star}\bigl(\widetilde V_c(w)-\widetilde V_{c^\star}(w)\bigr)>0.
\]
Alors il existe un voisinage \(U\) de \(w\) tel que tout \(z\in U\) admet le même contexte routé \(c^\star\).
\end{proposition}

\begin{proof}
Par continuité finie des fonctions \(\widetilde V_c\), on choisit \(U\) tel que \(|\widetilde V_c(z)-\widetilde V_c(w)|<\delta(w)/3\) pour tous \(c\) et tout \(z\in U\). L'argument standard de comparaison des branches donne alors la stabilité locale du contexte optimal.
\end{proof}

\begin{proposition}[Sensibilité continue aux paramètres]
Fixons \(w\in\R^d\). Pour \(\theta=(a,b)\in(0,\infty)\times[0,\infty)\), on définit
\[
J_{c,\theta}(u;w)
:=
\frac12\|u-w\|^2+a\,\dist(u,\Mcal_c)^2+b\,\|u\|^2,
\qquad
\widetilde V_{c,\theta}(w)
:=
\inf_{u\in\R^d}
\left\{
\frac{1}{1+2b}J_{c,\theta}(u;w)
\right\}.
\]
Alors, sur tout compact \(\Theta\subset(0,\infty)\times[0,\infty)\), l'application
\[
(w,\theta)\longmapsto \min_{c\in\Ccal}\widetilde V_{c,\theta}(w)
\]
est continue. Si, pour \((w_0,\theta_0)\), le minimiseur de \(c\mapsto \widetilde V_{c,\theta_0}(w_0)\) est unique, alors il reste stable par perturbation locale de \((w,\theta)\). En particulier, les frontières de changement de contexte se déplacent continûment lorsque les paramètres varient continûment.
\end{proposition}

\begin{proof}
Pour chaque \(c\), l'application \((w,\theta)\mapsto \widetilde V_{c,\theta}(w)\) est continue par les mêmes arguments de compacité et de convergence des minimiseurs que pour un problème proximal fortement convexe. Comme \(\Ccal\) est fini, le minimum fini de fonctions continues est continu. Ce point relève du théorème de Berge appliqué au cas discret et d'une paramétrisation fortement convexe \cite{RockafellarWets1998}.

Si le minimiseur en \((w_0,\theta_0)\) est unique, il existe un écart strict entre la meilleure valeur et la seconde meilleure. Par continuité, cet écart persiste dans un voisinage de \((w_0,\theta_0)\), ce qui stabilise localement le contexte optimal et déplace continûment les frontières de changement de contexte.
\end{proof}

% =================================================
\section{Routage souple comme relaxation entropique}
% =================================================

Le routage dur est discret. Pour obtenir une relaxation convexe, on passe au simplexe des contextes. Ce passage est une relaxation explicite du problème discret.

\begin{definition}[Simplexe des contextes]
On note
\[
\Delta_K
=
\left\{
p\in\R^K:\ p_c\ge 0,\ \sum_{c=1}^K p_c=1
\right\}.
\]
On interprète \(p_c\) comme le poids attribué au contexte \(c\), et on note
\[
\Delta_K^\circ
=
\{p\in\Delta_K:\ p_c>0\ \text{pour tout }c\}.
\]
\end{definition}

\begin{definition}[Fonctionnelle de routage souple]
Pour \(\beta>0\), on définit
\[
\Rcal_\beta(w,p)
=
\sum_{c=1}^K p_c\,\widetilde V_c(w)
+
\frac{1}{\beta}\sum_{c=1}^K p_c\log p_c,
\qquad p\in\Delta_K,
\]
avec la convention \(0\log 0=0\).
\end{definition}

\begin{proposition}[Existence et unicité du routage souple]
Pour tout \(w\), la fonctionnelle \(p\mapsto \Rcal_\beta(w,p)\) admet un unique minimiseur \(p^\star_\beta(w)\) dans \(\Delta_K\).
\end{proposition}

\begin{proof}
La fonction \(p\mapsto \sum_c p_c\log p_c\) est strictement convexe sur \(\Delta_K\) et \(p\mapsto \sum_c p_c\,\widetilde V_c(w)\) est affine. La somme est donc strictement convexe sur le compact \(\Delta_K\), d'où l'existence et l'unicité du minimiseur.

Montrons maintenant que le minimiseur appartient à l'intérieur \(\Delta_K^\circ\). Supposons au contraire que \(p^\star_i=0\) pour un certain indice \(i\). Comme \(\sum_c p^\star_c=1\), il existe un indice \(j\) tel que \(p^\star_j>0\). Pour \(t\in(0,p^\star_j)\), posons
\[
p(t)=p^\star+t(e_i-e_j)\in\Delta_K.
\]
Alors, à une constante près indépendante de \(t\),
\[
\phi(t)
=
t\,\widetilde V_i(w)
+
(p^\star_j-t)\,\widetilde V_j(w)
+
\frac{1}{\beta}\Bigl[t\log t+(p^\star_j-t)\log(p^\star_j-t)\Bigr].
\]
Donc
\[
\phi'(t)
=
\widetilde V_i(w)-\widetilde V_j(w)
+
\frac{1}{\beta}\Bigl(\log t-\log(p^\star_j-t)\Bigr).
\]
En particulier,
\[
\phi'(0^+)=-\infty
\qquad\text{car}\qquad
\log t\to-\infty\ \text{lorsque}\ t\downarrow0.
\]
Cela contredit la minimalité de \(p^\star\). Ainsi, tout minimiseur appartient à \(\Delta_K^\circ\), ce qui permet d'appliquer les conditions de Lagrange.
\end{proof}

\begin{theorem}[Forme explicite du routage souple]
Le minimiseur unique de \(\Rcal_\beta(w,\cdot)\) est donné par
\[
p_c^\star(w)
=
\frac{\exp\!\bigl(-\beta\,\widetilde V_c(w)\bigr)}
{\sum_{j=1}^K \exp\!\bigl(-\beta\,\widetilde V_j(w)\bigr)}.
\]
\end{theorem}

\begin{proof}
Comme le minimiseur appartient à \(\Delta_K^\circ\), on peut appliquer la méthode des multiplicateurs de Lagrange au problème
\[
\min_{p\in\Delta_K^\circ}\ \sum_{c=1}^K p_c\,\widetilde V_c(w)+\frac{1}{\beta}\sum_{c=1}^K p_c\log p_c.
\]
On obtient, pour chaque \(c\),
\[
\widetilde V_c(w)+\frac{1}{\beta}(1+\log p_c)+\lambda=0,
\]
donc
\[
p_c=\exp(-1-\beta\lambda)\,\exp\!\bigl(-\beta\widetilde V_c(w)\bigr).
\]
La contrainte \(\sum_c p_c=1\) donne la formule annoncée. Le résultat est la forme usuelle du softmin entropique \cite{BauschkeCombettes2017,Clarke1983}.
\end{proof}

\begin{corollary}[Valeur du problème de routage]
On a l'identité
\[
\inf_{p\in\Delta_K}\Rcal_\beta(w,p)
=
-\frac{1}{\beta}\log\!\left(\sum_{c=1}^K \exp\!\bigl(-\beta\,\widetilde V_c(w)\bigr)\right).
\]
\end{corollary}

\begin{proof}
C'est la valeur classique de l'entropie log-sum-exp obtenue en injectant la solution précédente dans \(\Rcal_\beta\).
\end{proof}

\begin{corollary}[Bornes de l'approximation entropique]
Pour tout \(w\),
\[
F(w)
\le
-\frac{1}{\beta}\log\!\left(\sum_{c=1}^K e^{-\beta \widetilde V_c(w)}\right)
\le
F(w)+\frac{\log K}{\beta}.
\]
\end{corollary}

\begin{proof}
La première inégalité est la propriété standard du \(\log\)-sum-exp comme lissage du minimum. Pour la seconde, on écrit
\[
\sum_{c=1}^K e^{-\beta \widetilde V_c(w)}
\le
K e^{-\beta F(w)},
\]
d'où
\[
-\frac{1}{\beta}\log\!\left(\sum_{c=1}^K e^{-\beta \widetilde V_c(w)}\right)
\ge
F(w)-\frac{\log K}{\beta},
\]
et la borne souhaitée en découle.
\end{proof}

\begin{corollary}[Routage dur comme limite]
Supposons que le minimiseur de \(c\mapsto \widetilde V_c(w)\) soit unique, noté \(c^\star(w)\). Alors
\[
p^\star_\beta(w)\xrightarrow[\beta\to+\infty]{} e_{c^\star(w)},
\]
où \(e_{c^\star}\) désigne le vecteur canonique.
\end{corollary}

\begin{proof}
Pour \(c\neq c^\star(w)\), on a
\[
p_c^\star(w)
=
\frac{\exp\!\bigl(-\beta(\widetilde V_c(w)-\widetilde V_{c^\star(w)}(w))\bigr)}
{1+\sum_{j\neq c^\star(w)}\exp\!\bigl(-\beta(\widetilde V_j(w)-\widetilde V_{c^\star(w)}(w))\bigr)}.
\]
Comme \(\widetilde V_c(w)-\widetilde V_{c^\star(w)}(w)>0\), le numérateur tend vers \(0\) exponentiellement lorsque \(\beta\to+\infty\). On obtient donc \(p_c^\star(w)\to0\) pour tout \(c\neq c^\star(w)\), et par la contrainte \(\sum_c p_c^\star(w)=1\), on en déduit \(p_{c^\star(w)}^\star(w)\to1\).
\end{proof}

\begin{remark}
Le routage souple est une relaxation entropique du problème discret. Il ne remplace pas le routage exact ; il l'approxime de manière stable.
\end{remark}

% =================================================
\section{Pré-roulage court et stabilisation}
% =================================================

Le routage exact peut être coûteux, car il exige la résolution complète de chaque problème contextuel. On introduit donc un certificat de calcul par pré-roulage court. Cette partie n'est pas un appendice : elle est reliée à la théorie de la valeur par une approximation consistante de \(\widetilde V_c\).

\begin{assumption}[Hypothèse de contraction]
Pour un contexte \(c\) fixé, on suppose que l'opérateur \(\Tcal_{c,\eta}(\cdot;w)\) est contractant sur \(\R^d\) avec facteur \(\rho_c\in(0,1)\), pour tout \(w\) dans le domaine considéré.
\end{assumption}

\begin{definition}[Pré-roulage court]
Pour un contexte \(c\), on pose
\[
u_c^{(0)}(w)=w,
\qquad
u_c^{(n+1)}(w)=\Tcal_{c,\eta}(u_c^{(n)}(w);w),
\qquad n\ge 0.
\]
\end{definition}

\begin{definition}[Valeur de pré-roulage]
Pour un entier \(N\ge 1\), on définit
\[
\widetilde V_c^{(N)}(w):=\widetilde J_c(u_c^{(N)}(w);w).
\]
\end{definition}

\begin{proposition}[Approximation de la valeur réduite]
Sous l'hypothèse de contraction ci-dessus, la suite \(u_c^{(N)}(w)\) converge vers \(u_c^\star(w)\). Par continuité de \(\widetilde J_c(\cdot;w)\), on a alors
\[
\widetilde V_c^{(N)}(w)\longrightarrow \widetilde V_c(w)
\quad\text{quand}\quad N\to\infty.
\]
\end{proposition}

\begin{proof}
La convergence des itérés découle du théorème du point fixe de Banach. La continuité de \(\widetilde J_c\) donne ensuite la convergence des valeurs.
\end{proof}

\begin{theorem}[Stabilisation du routage approché]
Supposons que le minimiseur de \(c\mapsto \widetilde V_c(w)\) soit unique, noté \(c^\star(w)\). Alors il existe \(N_0\ge 1\) tel que, pour tout \(N\ge N_0\),
\[
\hat c_N(w)=c^\star(w),
\]
où
\[
\hat c_N(w)\in\argmin_{c\in\Ccal}\widetilde V_c^{(N)}(w).
\]
\end{theorem}

\begin{proof}
Comme l'ensemble \(\Ccal\) est fini, \(\argmin_{c\in\Ccal}\widetilde V_c^{(N)}(w)\) est non vide pour tout \(N\). En outre, pour chaque \(c\),
\[
\widetilde V_c^{(N)}(w)\to \widetilde V_c(w).
\]
La convergence est donc uniforme sur \(\Ccal\). Posons
\[
\delta:=\min_{c\neq c^\star(w)}\bigl(\widetilde V_c(w)-\widetilde V_{c^\star(w)}(w)\bigr)>0.
\]
Il existe alors \(N_0\ge 1\) tel que, pour tout \(N\ge N_0\) et tout \(c\in\Ccal\),
\[
\bigl|\widetilde V_c^{(N)}(w)-\widetilde V_c(w)\bigr|<\frac{\delta}{3}.
\]
Pour tout \(c\neq c^\star(w)\), on obtient
\[
\widetilde V_c^{(N)}(w)
\ge
\widetilde V_c(w)-\frac{\delta}{3}
\ge
\widetilde V_{c^\star(w)}(w)+\delta-\frac{\delta}{3}
=
\widetilde V_{c^\star(w)}(w)+\frac{2\delta}{3},
\]
tandis que
\[
\widetilde V_{c^\star(w)}^{(N)}(w)
\le
\widetilde V_{c^\star(w)}(w)+\frac{\delta}{3}.
\]
Ainsi,
\[
\widetilde V_{c^\star(w)}^{(N)}(w)<\widetilde V_c^{(N)}(w)
\quad\text{pour tout } c\neq c^\star(w),
\]
ce qui prouve \(\hat c_N(w)=c^\star(w)\) pour tout \(N\ge N_0\).
\end{proof}

\begin{corollary}[Routage approché par certificat court]
On peut définir un routage approché par
\[
\hat c_N(w)\in\argmin_{c\in\Ccal}\widetilde V_c^{(N)}(w).
\]
Lorsque \(N\) est suffisamment grand, \(\hat c_N(w)\) coïncide avec \(\hat c(w)\) sous l'hypothèse d'unicité du minimiseur exact.
\end{corollary}

\begin{remark}
Le pré-roulage court n'est pas le principe théorique ; c'est un estimateur calculable de la valeur réduite. Le principe théorique reste la minimisation de \(\widetilde V_c(w)\) ou, de façon équivalente, la minimisation variationnelle sur le simplexe.
\end{remark}

% =================================================
\section{Interprétation variationnelle complète du routage}
% =================================================

Cette section rassemble la lecture unifiée du routage comme géométrie induite par une famille finie d'enveloppes de Moreau. C'est le point de vue conceptuel qui met le mieux en évidence l'objet original du manuscrit.

\begin{definition}[Enveloppe globale du routage]
On définit
\[
F(w):=\min_{c\in\Ccal}\widetilde V_c(w)
=
\min_{c\in\Ccal}e_{\mu_c g_c}(w).
\]
\end{definition}

\begin{proposition}[Enveloppe inférieure de plusieurs enveloppes de Moreau]
La fonction \(F\) est l'enveloppe inférieure d'une famille finie d'enveloppes de Moreau. En particulier :
\begin{enumerate}[leftmargin=2em]
\item \(F\) est continue et localement lipschitzienne ;
\item \(F\) est de classe \(C^1\) sur chaque cellule de dominance stricte ;
\item la strate de non-différentiabilité de \(F\) est contenue dans l'ensemble des points de bord \(\mathcal B\).
\end{enumerate}
\end{proposition}

\begin{proof}
Chaque \(\widetilde V_c=e_{\mu_c g_c}\) est \(C^1\) et localement lipschitzienne. Le minimum fini de fonctions continues est continu, et le minimum fini de fonctions localement lipschitziennes est localement lipschitzien. Si, en un point \(w\), un unique indice \(c^\star\) minimise \(c\mapsto \widetilde V_c(w)\), alors par continuité il reste minimiseur dans un voisinage de \(w\), et \(F\) y coïncide avec \(\widetilde V_{c^\star}\), donc est \(C^1\). Toute non-différentiabilité ne peut apparaître qu'en des points où au moins deux branches coïncident, c'est-à-dire dans \(\mathcal B\).
\end{proof}

\begin{proposition}[Sous-différentiel de la valeur globale]
Pour tout \(w\in\R^d\) et tout \(h\in\R^d\), la dérivée directionnelle de \(F\) existe et vérifie
\[
F'(w;h)=\min_{c\in\mathcal A(w)}\left\langle \frac{1}{\mu_c}\bigl(w-u_c^\star(w)\bigr),\,h\right\rangle.
\]
De plus, le sous-différentiel de Clarke de \(F\) au point \(w\) est
\[
\partial^C F(w)
=
\operatorname{conv}\left\{
\frac{1}{\mu_c}\bigl(w-u_c^\star(w)\bigr): c\in\mathcal A(w)
\right\}.
\]
\end{proposition}

\begin{proof}
C'est la règle standard du minimum fini de fonctions \(C^1\) localement lipschitziennes. Les branches actives sont exactement celles de \(\mathcal A(w)\), et les gradients locaux sont donnés par la formule de Moreau obtenue précédemment \cite{Clarke1983}.
\end{proof}

\begin{proposition}[Correspondance active globale]
La correspondance \(w\mapsto \mathcal A(w)\) a des valeurs non vides et compactes, est semi-continue supérieurement sur \(\R^d\), et est localement constante sur chaque cellule de dominance stricte.
\end{proposition}

\begin{proof}
Les valeurs sont des sous-ensembles non vides d'un ensemble fini, donc compactes. La semi-continuité supérieure a déjà été établie via la fermeture du graphe. La constance locale sur une cellule stricte résulte de l'existence d'un écart strict entre le contexte actif et tous les autres.
\end{proof}

\begin{definition}[Frontières de changement de contexte]
Pour chaque paire \(c\neq j\), on pose
\[
\Gamma_{c,j}(w):=\widetilde V_c(w)-\widetilde V_j(w),
\qquad
\Sigma_{c,j}:=\{w\in\R^d:\Gamma_{c,j}(w)=0\}.
\]
Les ensembles \(\Sigma_{c,j}\) sont les frontières paires de changement de contexte.
\end{definition}

\begin{proposition}[Stabilité locale des frontières]
Pour chaque paire \(c\neq j\), l'ensemble \(\Sigma_{c,j}\) est fermé. Si \(w_0\in\Sigma_{c,j}\) et si
\[
\nabla \Gamma_{c,j}(w_0)\neq 0,
\]
alors il existe un voisinage \(U\) de \(w_0\) dans lequel \(\Sigma_{c,j}\cap U\) est une hypersurface \(C^1\) de \(\R^d\).
\end{proposition}

\begin{proof}
La fermeture de \(\Sigma_{c,j}\) découle de la continuité de \(\Gamma_{c,j}\). Si \(\nabla \Gamma_{c,j}(w_0)\neq 0\), le théorème des fonctions implicites s'applique à l'équation \(\Gamma_{c,j}(w)=0\) au voisinage de \(w_0\), ce qui fournit une hypersurface \(C^1\) locale.
\end{proof}

\begin{proposition}[Stabilité sous marge stricte]
Supposons qu'en un point \(w\in\R^d\), le minimiseur \(c^\star\) de \(c\mapsto \widetilde V_c(w)\) soit unique et qu'il existe une marge
\[
\delta(w):=\min_{c\neq c^\star}\bigl(\widetilde V_c(w)-\widetilde V_{c^\star}(w)\bigr)>0.
\]
Alors il existe un voisinage \(U\) de \(w\) tel que tout \(z\in U\) admet le même contexte routé \(c^\star\).
\end{proposition}

\begin{proof}
Par continuité finie des fonctions \(\widetilde V_c\), on choisit \(U\) tel que \(|\widetilde V_c(z)-\widetilde V_c(w)|<\delta(w)/3\) pour tous \(c\) et tout \(z\in U\). L'argument standard de comparaison des branches donne alors la stabilité locale du contexte optimal.
\end{proof}

\begin{remark}
Cette lecture complète montre que le routage est une géométrie de choix entre branches convexes : le noyau PCH fournit les branches, la valeur réduite les compare, la valeur globale en est l'enveloppe inférieure, et les frontières sont les lieux où cette enveloppe change de branche.
\end{remark}

% =================================================
\section{Programmation disjonctive et mélange d'experts variationnel}
% =================================================

Le problème global de routage peut être lu comme un programme disjonctif fini :
\[
F(w)=\min_{c\in\Ccal}\widetilde V_c(w)
=
\min_{c\in\Ccal}\inf_{u\in\R^d}\widetilde J_c(u;w).
\]
Chaque contexte \(c\) définit une branche convexe du problème, et le routage dur choisit exactement une branche active.

\begin{definition}[Banque d'experts variationnels]
Pour chaque contexte \(c\), on appelle expert variationnel l'application
\[
E_c:\R^d\to\R^d,\qquad E_c(w):=u_c^\star(w),
\]
et on appelle gate variationnel la sélection
\[
G(w)\in\argmin_{c\in\Ccal}\widetilde V_c(w).
\]
Dans la version souple, le gate devient la loi de Gibbs
\[
G_\beta(w)=p^\star_\beta(w).
\]
\end{definition}

\begin{proposition}[Lecture disjonctive du routage]
Le routage dur correspond à la sélection d'une branche dans un programme disjonctif fini. Le routage souple correspond à une régularisation entropique de cette sélection, dont la solution est la loi de Gibbs.
\end{proposition}

\begin{proof}
La première assertion découle directement de l'écriture
\[
F(w)=\min_{c\in\Ccal}\widetilde V_c(w),
\]
où chaque \(\widetilde V_c\) est la valeur d'un sous-problème convexe. La seconde assertion découle de la formule exacte du routage souple :
\[
p^\star_\beta(w)
=
\argmin_{p\in\Delta_K}
\left\{
\sum_{c=1}^K p_c\,\widetilde V_c(w)
+
\frac{1}{\beta}\sum_{c=1}^K p_c\log p_c
\right\}.
\]
Le premier terme sélectionne une branche, le second la lisse par entropie.
\end{proof}

\begin{remark}
Cette lecture est l'analogue variationnel d'un mélange d'experts, mais avec une différence essentielle : le gate n'est pas appris comme un réseau séparé, il est défini par la résolution exacte d'un problème d'optimisation.
\end{remark}

% =================================================
\section{Lien avec les heuristiques de routage et la littérature}
% =================================================

Les heuristiques de routage utilisées en pratique peuvent être vues comme des approximations de la valeur réduite ou de son gain normalisé. Le tableau ci-dessous synthétise leur lecture variationnelle.

\begin{center}
\begin{tabular}{@{}lll@{}}
\toprule
Heuristique & Quantité approchée & Rôle \\ \midrule
Distance de projection & \(\widetilde V_c(w)\) & estimateur local \\
Gain de stabilité & \(\widetilde\Gamma_c(w)\) & diagnostic de comparaison \\
Pré-roulage court & \(\widetilde V_c^{(N)}(w)\) & approximation calculable \\
Routage appris & \(\argmin_{c\in\Ccal}\widetilde V_c(w)\) & approximation supervisée \\ \bottomrule
\end{tabular}
\end{center}

\begin{proposition}[Lecture théorique des heuristiques]
Les heuristiques de routage peuvent être interprétées comme des estimateurs de \(\widetilde V_c(w)\), de \(\widetilde\Gamma_c(w)\), ou de leur minimiseur. Leur qualité dépend de leur fidélité à la structure variationnelle du problème.
\end{proposition}

\begin{proof}
Chaque heuristique cherche à approximer une quantité qui sélectionne le contexte minimisant le critère adopté. Le cadre ci-dessus fournit précisément cette quantité : une valeur réduite normalisée, sa version entropiquement lissée, ou son approximation par pré-roulage court.
\end{proof}

\begin{remark}
Le cadre ici présenté se place naturellement à l'interface de l'analyse convexe, de la théorie des fonctions valeur, de la programmation disjonctive et des modèles de type \emph{mixture of experts}. Les références classiques pertinentes sont celles de Moreau, Rockafellar, Berge, Clarke, ainsi que Bauschke--Combettes \cite{Moreau1965,Rockafellar1970,Berge1963,Clarke1983,BauschkeCombettes2017,RockafellarWets1998}.
\end{remark}

% =================================================
\section{Conclusion}
% =================================================

Le noyau PCH n'est pas le problème à résoudre ; il est déjà la brique mathématique centrale. Le vrai problème est le routage : comment choisir le contexte de manière rigoureuse, sans introduire un contournement externe.

La réponse proposée ici est écrite de façon exacte : le routage est un problème d'optimisation sur l'espace produit \(\R^d\times\Ccal\), et sa relaxation entropique sur le simplexe fournit une version souple. La formulation canonique est
\[
(\hat u(w),\hat c(w))\in\argmin_{(u,c)\in\R^d\times\Ccal}\mathcal E_{\mathrm{route}}(w;u,c),
\]
équivalente à
\[
\hat c(w)\in\argmin_{c\in\Ccal}\widetilde V_c(w),
\qquad
\widetilde V_c(w)=e_{\mu_c g_c}(w).
\]
La version souple est donnée par
\[
p^\star_\beta(w)
=
\argmin_{p\in\Delta_K}\Rcal_\beta(w,p).
\]

Cette formulation conserve strictement le noyau PCH, définit explicitement le critère de sélection du contexte, donne une écriture exacte du routage comme problème sur \(\R^d\times\Ccal\), admet une relaxation convexe sur le simplexe, admet un certificat de pré-roulage court comme approximation, et donne une interprétation variationnelle propre du choix de contexte.

Le routage s'inscrit ainsi dans la formulation mathématique du cadre proposé.

\begin{thebibliography}{99}

\bibitem{Moreau1965}
J.-J. Moreau,
\newblock Proximité et dualité dans un espace hilbertien.
\newblock \emph{Bull. Soc. Math. France}, 93:273--299, 1965.

\bibitem{Rockafellar1970}
R. T. Rockafellar,
\newblock \emph{Convex Analysis}.
\newblock Princeton University Press, 1970.

\bibitem{Berge1963}
C. Berge,
\newblock \emph{Topological Spaces: Including a Treatment of Multi-Valued Functions, Vector Spaces and Convexity}.
\newblock Oliver \& Boyd, 1963.

\bibitem{Clarke1983}
F. H. Clarke,
\newblock \emph{Optimization and Nonsmooth Analysis}.
\newblock Wiley, 1983.

\bibitem{BauschkeCombettes2017}
H. H. Bauschke and P. L. Combettes,
\newblock \emph{Convex Analysis and Monotone Operator Theory in Hilbert Spaces}.
\newblock 2nd ed., Springer, 2017.

\bibitem{RockafellarWets1998}
R. T. Rockafellar and R. J.-B. Wets,
\newblock \emph{Variational Analysis}.
\newblock Springer, 1998.

\end{thebibliography}

\end{document}
